
% =====================================================
% FICHE 1 - Chaîne radar instable
% =====================================================
\begin{center}
{\LARGE \textbf{FICHE DE RISQUE}}\\[6pt]
{\large Ref : \textbf{R001} \hfill Lot : \underline{Chaîne radar instable}}\\[6pt]
\hrule
\end{center}

\renewcommand{\arraystretch}{1.2}
\begin{tabular}{|p{5.2cm}|p{10.8cm}|}
\hline
\textbf{Libellé du risque} & Chaîne radar instable \\ \hline
\textbf{Classe (acteur exposé)} & Équipe globale \\ \hline
\textbf{Responsable} & Référent RF / Chef de projet \\ \hline
\textbf{Causes potentielles ou réelles} &
\begin{minipage}[t]{\linewidth}\vspace{2pt}
\begin{itemize}\setlength\itemsep{2pt}
\item Dérive de fréquence, bruit RF ou mauvais étalonnage
\item Mauvaise isolation entre émetteur et récepteur
\item Composants non adaptés ou vieillissants
\end{itemize}\vspace{2pt}
\end{minipage} \\ \hline
\textbf{Conséquences / Impacts possibles} &
\begin{minipage}[t]{\linewidth}\vspace{2pt}
\begin{itemize}\setlength\itemsep{2pt}
\item Données I/Q inutilisables
\item Retard de validation du prototype
\item Reprise complète du montage ou recalibrage
\end{itemize}\vspace{2pt}
\end{minipage} \\ \hline
\textbf{Actions de réduction / prévention} &
\begin{minipage}[t]{\linewidth}\vspace{2pt}
\begin{itemize}\setlength\itemsep{2pt}
\item Effectuer un étalonnage initial complet
\item Ajouter des étapes de test RF intermédiaires
\item Documenter les configurations matérielles
\end{itemize}\vspace{2pt}
\end{minipage} \\ \hline
\textbf{Probabilité (P)} & 3 \\ \hline
\textbf{Gravité (Gc,Gd,Gp)} & Gc = 0, Gd = 4, Gp = 4 \\ \hline
\textbf{Coefficient de résolution (R)} & 4 \\ \hline
\textbf{Sévérité globale (SG)} & 96 \\ \hline
\textbf{Niveau d'exposition} & Critique \\ \hline
\end{tabular}
\newpage
% =====================================================
% FICHE 2 - Perte ou corruption de données expérimentales
% =====================================================
\begin{center}
{\LARGE \textbf{FICHE DE RISQUE}}\\[6pt]
{\large Ref : \textbf{R002} \hfill Lot : \underline{Perte ou corruption de données expérimentales}}\\[6pt]
\hrule
\end{center}

\begin{tabular}{|p{5.2cm}|p{10.8cm}|}
\hline
\textbf{Libellé du risque} & Perte ou corruption de données expérimentales \\ \hline
\textbf{Classe (acteur exposé)} & Équipe traitement \\ \hline
\textbf{Responsable} & Responsable traitement SAR \\ \hline
\textbf{Causes potentielles ou réelles} &
\begin{minipage}[t]{\linewidth}\vspace{2pt}
\begin{itemize}\setlength\itemsep{2pt}
\item Erreur de stockage ou panne SSD
\item Mauvaise sauvegarde des acquisitions I/Q
\item Mauvais versioning
\end{itemize}\vspace{2pt}
\end{minipage} \\ \hline
\textbf{Conséquences / Impacts possibles} &
\begin{minipage}[t]{\linewidth}\vspace{2pt}
\begin{itemize}\setlength\itemsep{2pt}
\item Données SAR inutilisables
\item Reprise complète des acquisitions
\item Retard dans les tests et l’analyse
\end{itemize}\vspace{2pt}
\end{minipage} \\ \hline
\textbf{Actions de réduction / prévention} &
\begin{minipage}[t]{\linewidth}\vspace{2pt}
\begin{itemize}\setlength\itemsep{2pt}
\item Sauvegarde automatique sur serveur IMT
\item Double stockage local + cloud
\item Versioning des fichiers
\end{itemize}\vspace{2pt}
\end{minipage} \\ \hline
\textbf{Probabilité (P)} & 2 \\ \hline
\textbf{Gravité (Gc,Gd,Gp)} & Gc = 0, Gd = 3, Gp = 3 \\ \hline
\textbf{Coefficient de résolution (R)} & 4 \\ \hline
\textbf{Sévérité globale (SG)} & 48 \\ \hline
\textbf{Niveau d'exposition} & Modéré \\ \hline
\end{tabular}
\newpage
% =====================================================
% FICHE 3 - Problème d’expérimentation matérielle
% =====================================================
\begin{center}
{\LARGE \textbf{FICHE DE RISQUE}}\\[6pt]
{\large Ref : \textbf{R003} \hfill Lot : \underline{Problème d’expérimentation matérielle}}\\[6pt]
\hrule
\end{center}

\begin{tabular}{|p{5.2cm}|p{10.8cm}|}
\hline
\textbf{Libellé du risque} & Problème d’expérimentation (matérielle) \\ \hline
\textbf{Classe (acteur exposé)} & Équipe hardware \\ \hline
\textbf{Responsable} & Référent intégration matérielle \\ \hline
\textbf{Causes potentielles ou réelles} &
\begin{minipage}[t]{\linewidth}\vspace{2pt}
\begin{itemize}\setlength\itemsep{2pt}
\item Alimentation insuffisante ou défaillante
\item Connectique incompatible entre modules
\item Mauvais câblage ou surchauffe
\end{itemize}\vspace{2pt}
\end{minipage} \\ \hline
\textbf{Conséquences / Impacts possibles} &
\begin{minipage}[t]{\linewidth}\vspace{2pt}
\begin{itemize}\setlength\itemsep{2pt}
\item Perte de temps sur les essais
\item Arrêt du développement
\item Reconfiguration matérielle non prévue
\end{itemize}\vspace{2pt}
\end{minipage} \\ \hline
\textbf{Actions de réduction / prévention} &
\begin{minipage}[t]{\linewidth}\vspace{2pt}
\begin{itemize}\setlength\itemsep{2pt}
\item Prévoir des composants de rechange
\item Réaliser des tests unitaires avant chaque essai
\item Contrôler thermiquement le montage
\end{itemize}\vspace{2pt}
\end{minipage} \\ \hline
\textbf{Probabilité (P)} & 2 \\ \hline
\textbf{Gravité (Gc,Gd,Gp)} & Gc = 4, Gd = 3, Gp = 4 \\ \hline
\textbf{Coefficient de résolution (R)} & 2 \\ \hline
\textbf{Sévérité globale (SG)} & 44 \\ \hline
\textbf{Niveau d'exposition} & Modéré \\ \hline
\end{tabular}

\newpage
% =====================================================
% FICHE 4 - Fusion multi-capteurs
% =====================================================
\begin{center}
{\LARGE \textbf{FICHE DE RISQUE}}\\[6pt]
{\large Ref : \textbf{R004} \hfill Lot : \underline{Fusion multi-capteurs complexe ou imprécise}}\\[6pt]
\hrule
\end{center}

\begin{tabular}{|p{5.2cm}|p{10.8cm}|}
\hline
\textbf{Libellé du risque} & Fusion multi-capteurs complexe ou imprécise \\ \hline
\textbf{Classe (acteur exposé)} & Équipe algorithmie / traitement SAR \\ \hline
\textbf{Responsable} & Responsable synchronisation / traitement \\ \hline
\textbf{Causes potentielles ou réelles} &
\begin{minipage}[t]{\linewidth}\vspace{2pt}
\begin{itemize}\setlength\itemsep{2pt}
\item Désynchronisation temporelle entre radars
\item Mauvaise calibration spatiale
\item Interférences CEM entre capteurs
\end{itemize}\vspace{2pt}
\end{minipage} \\ \hline
\textbf{Conséquences / Impacts possibles} &
\begin{minipage}[t]{\linewidth}\vspace{2pt}
\begin{itemize}\setlength\itemsep{2pt}
\item Images SAR floues ou incohérentes
\item Résultats scientifiques invalides
\item Difficulté à prouver la faisabilité multi-capteurs
\end{itemize}\vspace{2pt}
\end{minipage} \\ \hline
\textbf{Actions de réduction / prévention} &
\begin{minipage}[t]{\linewidth}\vspace{2pt}
\begin{itemize}\setlength\itemsep{2pt}
\item Synchroniser les acquisitions (horloge commune)
\item Tester à 2 capteurs avant extension
\item Évaluer la CEM en amont
\end{itemize}\vspace{2pt}
\end{minipage} \\ \hline
\textbf{Probabilité (P)} & 2 \\ \hline
\textbf{Gravité (Gc,Gd,Gp)} & Gc = 0, Gd = 3, Gp = 4 \\ \hline
\textbf{Coefficient de résolution (R)} & 3 \\ \hline
\textbf{Sévérité globale (SG)} & 42 \\ \hline
\textbf{Niveau d'exposition} & Modéré \\ \hline
\end{tabular}

\newpage
% =====================================================
% FICHE 5 - Non-conformité réglementaire
% =====================================================
\begin{center}
{\LARGE \textbf{FICHE DE RISQUE}}\\[6pt]
{\large Ref : \textbf{R005} \hfill Lot : \underline{Non-conformité réglementaire (RF / sécurité)}}\\[6pt]
\hrule
\end{center}

\begin{tabular}{|p{5.2cm}|p{10.8cm}|}
\hline
\textbf{Libellé du risque} & Non-conformité réglementaire (RF / sécurité) \\ \hline
\textbf{Classe (acteur exposé)} & Encadrants responsables légaux \\ \hline
\textbf{Responsable} & Encadrant RF / Responsable sécurité IMT \\ \hline
\textbf{Causes potentielles ou réelles} &
\begin{minipage}[t]{\linewidth}\vspace{2pt}
\begin{itemize}\setlength\itemsep{2pt}
\item Utilisation de bandes de fréquences non autorisées
\item Essais non déclarés ou hors site IMT
\item Absence de validation sécurité
\end{itemize}\vspace{2pt}
\end{minipage} \\ \hline
\textbf{Conséquences / Impacts possibles} &
\begin{minipage}[t]{\linewidth}\vspace{2pt}
\begin{itemize}\setlength\itemsep{2pt}
\item Interdiction d’expérimentation
\item Retard du projet
\item Image négative du groupe
\end{itemize}\vspace{2pt}
\end{minipage} \\ \hline
\textbf{Actions de réduction / prévention} &
\begin{minipage}[t]{\linewidth}\vspace{2pt}
\begin{itemize}\setlength\itemsep{2pt}
\item Vérifier les fréquences avant test
\item Réaliser essais en environnement contrôlé
\item Validation sécurité par encadrant
\end{itemize}\vspace{2pt}
\end{minipage} \\ \hline
\textbf{Probabilité (P)} & 1 \\ \hline
\textbf{Gravité (Gc,Gd,Gp)} & Gc = 4, Gd = 4, Gp = 4 \\ \hline
\textbf{Coefficient de résolution (R)} & 2 \\ \hline
\textbf{Sévérité globale (SG)} & 24 \\ \hline
\textbf{Niveau d'exposition} & Modéré \\ \hline
\end{tabular}

\newpage
% =====================================================
% FICHE 6 - Démotivation ou perte de cohésion
% =====================================================
\begin{center}
{\LARGE \textbf{FICHE DE RISQUE}}\\[6pt]
{\large Ref : \textbf{R006} \hfill Lot : \underline{Démotivation ou perte de cohésion d’équipe}}\\[6pt]
\hrule
\end{center}

\begin{tabular}{|p{5.2cm}|p{10.8cm}|}
\hline
\textbf{Libellé du risque} & Démotivation ou perte de cohésion d’équipe \\ \hline
\textbf{Classe (acteur exposé)} & Équipe globale \\ \hline
\textbf{Responsable} & Chef de projet / Référent humain \\ \hline
\textbf{Causes potentielles ou réelles} &
\begin{minipage}[t]{\linewidth}\vspace{2pt}
\begin{itemize}\setlength\itemsep{2pt}
\item Fatigue, déséquilibre des contributions
\item Difficulté technique prolongée sans résultat
\item Mauvais partage des réussites
\end{itemize}\vspace{2pt}
\end{minipage} \\ \hline
\textbf{Conséquences / Impacts possibles} &
\begin{minipage}[t]{\linewidth}\vspace{2pt}
\begin{itemize}\setlength\itemsep{2pt}
\item Baisse de productivité
\item Climat d’équipe détérioré
\item Risque sur la qualité des livrables
\end{itemize}\vspace{2pt}
\end{minipage} \\ \hline
\textbf{Actions de réduction / prévention} &
\begin{minipage}[t]{\linewidth}\vspace{2pt}
\begin{itemize}\setlength\itemsep{2pt}
\item Points d’équipe réguliers
\item Répartition équitable des tâches
\item Valoriser les progrès individuels
\end{itemize}\vspace{2pt}
\end{minipage} \\ \hline
\textbf{Probabilité (P)} & 1 \\ \hline
\textbf{Gravité (Gc,Gd,Gp)} & Gc = 0, Gd = 3, Gp = 4 \\ \hline
\textbf{Coefficient de résolution (R)} & 3 \\ \hline
\textbf{Sévérité globale (SG)} & 21 \\ \hline
\textbf{Niveau d'exposition} & Acceptable \\ \hline
\end{tabular}

\newpage
% =====================================================
% FICHE 7 - Retard sur un jalon
% =====================================================
\begin{center}
{\LARGE \textbf{FICHE DE RISQUE}}\\[6pt]
{\large Ref : \textbf{R007} \hfill Lot : \underline{Retard sur un jalon}}\\[6pt]
\hrule
\end{center}

\begin{tabular}{|p{5.2cm}|p{10.8cm}|}
\hline
\textbf{Libellé du risque} & Retard sur un jalon \\ \hline
\textbf{Classe (acteur exposé)} & Planning global \\ \hline
\textbf{Responsable} & Chef de projet \\ \hline
\textbf{Causes potentielles ou réelles} &
\begin{minipage}[t]{\linewidth}\vspace{2pt}
\begin{itemize}\setlength\itemsep{2pt}
\item Complexité technique sous-estimée
\item Indisponibilité d’un membre clé
\item Réunion de validation reportée
\end{itemize}\vspace{2pt}
\end{minipage} \\ \hline
\textbf{Conséquences / Impacts possibles} &
\begin{minipage}[t]{\linewidth}\vspace{2pt}
\begin{itemize}\setlength\itemsep{2pt}
\item Décalage du calendrier global
\item Réduction du périmètre des tests
\item Pénalité sur la qualité du livrable ProCom
\end{itemize}\vspace{2pt}
\end{minipage} \\ \hline
\textbf{Actions de réduction / prévention} &
\begin{minipage}[t]{\linewidth}\vspace{2pt}
\begin{itemize}\setlength\itemsep{2pt}
\item Suivi hebdomadaire
\item Validation intermédiaire à chaque jalon
\item Redéfinir les priorités en cas de retard
\end{itemize}\vspace{2pt}
\end{minipage} \\ \hline
\textbf{Probabilité (P)} & 2 \\ \hline
\textbf{Gravité (Gc,Gd,Gp)} & Gc = 0, Gd = 4, Gp = 0 \\ \hline
\textbf{Coefficient de résolution (R)} & 2 \\ \hline
\textbf{Sévérité globale (SG)} & 16 \\ \hline
\textbf{Niveau d'exposition} & Acceptable \\ \hline
\end{tabular}

\newpage
% =====================================================
% FICHE 8 - Surcharge de calcul / performance logicielle
% =====================================================
\begin{center}
{\LARGE \textbf{FICHE DE RISQUE}}\\[6pt]
{\large Ref : \textbf{R008} \hfill Lot : \underline{Surcharge de calcul / performance logicielle}}\\[6pt]
\hrule
\end{center}

\begin{tabular}{|p{5.2cm}|p{10.8cm}|}
\hline
\textbf{Libellé du risque} & Surcharge de calcul / contraintes de performance logicielle \\ \hline
\textbf{Classe (acteur exposé)} & Équipe traitement SAR / simulation \\ \hline
\textbf{Responsable} & Responsable algorithmie / simulation \\ \hline
\textbf{Causes potentielles ou réelles} &
\begin{minipage}[t]{\linewidth}\vspace{2pt}
\begin{itemize}\setlength\itemsep{2pt}
\item Scripts non optimisés
\item Ressources serveur insuffisantes
\item Traitement SAR trop long
\end{itemize}\vspace{2pt}
\end{minipage} \\ \hline
\textbf{Conséquences / Impacts possibles} &
\begin{minipage}[t]{\linewidth}\vspace{2pt}
\begin{itemize}\setlength\itemsep{2pt}
\item Résultats non exploitables en temps voulu
\item Blocage du pipeline de traitement
\item Retard sur les livrables
\end{itemize}\vspace{2pt}
\end{minipage} \\ \hline
\textbf{Actions de réduction / prévention} &
\begin{minipage}[t]{\linewidth}\vspace{2pt}
\begin{itemize}\setlength\itemsep{2pt}
\item Optimiser les scripts Python/C++
\item Réserver les serveurs IMT à l’avance
\item Réduire la taille des jeux de données test
\end{itemize}\vspace{2pt}
\end{minipage} \\ \hline
\textbf{Probabilité (P)} & 2 \\ \hline
\textbf{Gravité (Gc,Gd,Gp)} & Gc = 0, Gd = 4, Gp = 0 \\ \hline
\textbf{Coefficient de résolution (R)} & 2 \\ \hline
\textbf{Sévérité globale (SG)} & 16 \\ \hline
\textbf{Niveau d'exposition} & Acceptable \\ \hline
\end{tabular}

\newpage
% =====================================================
% FICHE 9 - Dépassement du budget matériel
% =====================================================
\begin{center}
{\LARGE \textbf{FICHE DE RISQUE}}\\[6pt]
{\large Ref : \textbf{R009} \hfill Lot : \underline{Dépassement du budget matériel}}\\[6pt]
\hrule
\end{center}

\begin{tabular}{|p{5.2cm}|p{10.8cm}|}
\hline
\textbf{Libellé du risque} & Dépassement du budget matériel \\ \hline
\textbf{Classe (acteur exposé)} & Équipe globale \\ \hline
\textbf{Responsable} & Chef de projet / Encadrant IMT \\ \hline
\textbf{Causes potentielles ou réelles} &
\begin{minipage}[t]{\linewidth}\vspace{2pt}
\begin{itemize}\setlength\itemsep{2pt}
\item Besoins matériels sous-évalués
\item Erreurs de commande ou achats multiples
\item Coûts non anticipés
\end{itemize}\vspace{2pt}
\end{minipage} \\ \hline
\textbf{Conséquences / Impacts possibles} &
\begin{minipage}[t]{\linewidth}\vspace{2pt}
\begin{itemize}\setlength\itemsep{2pt}
\item Impossibilité d’acquérir du matériel critique
\item Nécessité de revoir les priorités d’essai
\end{itemize}\vspace{2pt}
\end{minipage} \\ \hline
\textbf{Actions de réduction / prévention} &
\begin{minipage}[t]{\linewidth}\vspace{2pt}
\begin{itemize}\setlength\itemsep{2pt}
\item Valider les dépenses avec l’équipe
\item Mutualiser le matériel IMT
\item Tenir un tableau de suivi budgétaire
\end{itemize}\vspace{2pt}
\end{minipage} \\ \hline
\textbf{Probabilité (P)} & 1 \\ \hline
\textbf{Gravité (Gc,Gd,Gp)} & Gc = 4, Gd = 0, Gp = 0 \\ \hline
\textbf{Coefficient de résolution (R)} & 3 \\ \hline
\textbf{Sévérité globale (SG)} & 12 \\ \hline
\textbf{Niveau d'exposition} & Acceptable \\ \hline
\end{tabular}

\newpage
% =====================================================
% FICHE 10 - Répartition inégale des tâches
% =====================================================
\begin{center}
{\LARGE \textbf{FICHE DE RISQUE}}\\[6pt]
{\large Ref : \textbf{R010} \hfill Lot : \underline{Répartition inégale des tâches / absences}}\\[6pt]
\hrule
\end{center}

\begin{tabular}{|p{5.2cm}|p{10.8cm}|}
\hline
\textbf{Libellé du risque} & Répartition inégale des tâches / absences \\ \hline
\textbf{Classe (acteur exposé)} & Équipe globale \\ \hline
\textbf{Responsable} & Chef de projet \\ \hline
\textbf{Causes potentielles ou réelles} &
\begin{minipage}[t]{\linewidth}\vspace{2pt}
\begin{itemize}\setlength\itemsep{2pt}
\item Mauvaise communication interne
\item Manque de coordination ou motivation
\item Absence d’un membre sur une période critique
\end{itemize}\vspace{2pt}
\end{minipage} \\ \hline
\textbf{Conséquences / Impacts possibles} &
\begin{minipage}[t]{\linewidth}\vspace{2pt}
\begin{itemize}\setlength\itemsep{2pt}
\item Retard sur certaines tâches
\item Déséquilibre dans la charge de travail
\item Diminution de la qualité du livrable
\end{itemize}\vspace{2pt}
\end{minipage} \\ \hline
\textbf{Actions de réduction / prévention} &
\begin{minipage}[t]{\linewidth}\vspace{2pt}
\begin{itemize}\setlength\itemsep{2pt}
\item Répartition claire des livrables
\item Suivi par binômes
\item Plan de secours en cas d’absence
\end{itemize}\vspace{2pt}
\end{minipage} \\ \hline
\textbf{Probabilité (P)} & 2 \\ \hline
\textbf{Gravité (Gc,Gd,Gp)} & Gc = 0, Gd = 2, Gp = 0 \\ \hline
\textbf{Coefficient de résolution (R)} & 2 \\ \hline
\textbf{Sévérité globale (SG)} & 8 \\ \hline
\textbf{Niveau d'exposition} & Acceptable \\ \hline
\end{tabular}

\newpage
% =====================================================
% FICHE 11 - Défaut de communication
% =====================================================
\begin{center}
{\LARGE \textbf{FICHE DE RISQUE}}\\[6pt]
{\large Ref : \textbf{R011} \hfill Lot : \underline{Défaut de communication inter-équipe ou encadrants}}\\[6pt]
\hrule
\end{center}

\begin{tabular}{|p{5.2cm}|p{10.8cm}|}
\hline
\textbf{Libellé du risque} & Défaut de communication inter-équipe ou avec encadrants \\ \hline
\textbf{Classe (acteur exposé)} & Équipe projet + encadrants \\ \hline
\textbf{Responsable} & Référent communication / Chef de projet \\ \hline
\textbf{Causes potentielles ou réelles} &
\begin{minipage}[t]{\linewidth}\vspace{2pt}
\begin{itemize}\setlength\itemsep{2pt}
\item Absence de compte rendu régulier
\item Manque de clarté sur les responsabilités
\item Multiplicité des canaux d’échange
\end{itemize}\vspace{2pt}
\end{minipage} \\ \hline
\textbf{Conséquences / Impacts possibles} &
\begin{minipage}[t]{\linewidth}\vspace{2pt}
\begin{itemize}\setlength\itemsep{2pt}
\item Mauvaise cohérence des sous-parties
\item Décisions non validées ou contradictoires
\item Risque de réorientation tardive du projet
\end{itemize}\vspace{2pt}
\end{minipage} \\ \hline
\textbf{Actions de réduction / prévention} &
\begin{minipage}[t]{\linewidth}\vspace{2pt}
\begin{itemize}\setlength\itemsep{2pt}
\item Réunion hebdomadaire obligatoire
\item Compte rendu partagé sur Drive
\item Référent communication désigné
\end{itemize}\vspace{2pt}
\end{minipage} \\ \hline
\textbf{Probabilité (P)} & 1 \\ \hline
\textbf{Gravité (Gc,Gd,Gp)} & Gc = 0, Gd = 3, Gp = 1 \\ \hline
\textbf{Coefficient de résolution (R)} & 1 \\ \hline
\textbf{Sévérité globale (SG)} & 4 \\ \hline
\textbf{Niveau d'exposition} & Acceptable \\ \hline
\end{tabular}
